% =========================================================
\chapter{Lecture 9 : Adversarial thinking II}
% =========================================================

\begin{rembox}[title=In a nutshell]
\textbf{What is this lecture about?} Lecture 8 asked ``what can go wrong?''. This one asks
\textbf{``what not to do?''}: the \textbf{Common Weakness Enumeration (CWE)}, MITRE's database of
the software errors that most often lead to vulnerabilities, grouped in three classes:
\textbf{insecure interaction between components}, \textbf{risky resource management} and
\textbf{porous defenses}. Most examples are on the web, so the lecture starts with a short
\textbf{web cheat sheet} (HTTP, cookies, GET and POST, HTML, PHP). We then dissect three classic
injection-style attacks: \textbf{OS command injection}, \textbf{cross-site scripting (XSS)} and
\textbf{cross-site request forgery (CSRF)}, with the \textbf{Same Origin Policy} and the defenses
against CSRF. \\
\textbf{Why is it useful?} XSS, SQL injection, CSRF and command injection sit in the top 10 of the
most dangerous software errors year after year. They all come from one mistake, using input you do
not control without sanitizing it, and one mental model, the confused deputy, explains CSRF.
\end{rembox}

\section{Web preliminaries}

\begin{rembox}[title=Why this cheat sheet]
Not a guide to web programming: only the basic concepts needed to understand the security problems
of the lecture. Up to now the course had two scenarios: a user acting \textbf{locally} on their
own device (local authentication, access control and execution), and a user acting
\textbf{remotely} on a server (the device is merely a screen). On the web, \textbf{browser and
server collaborate}: remote authentication, remote access control, and \textbf{mixed program
execution} (code runs on both sides). Two tools: \textbf{HTTP}, the protocol between browsers and
servers, and \textbf{PHP}, a language embedding dynamic, server-side elements in web pages.
\end{rembox}

\begin{defnbox}[title={Where we are: the top layers of the Internet (OSI model)}]
\begin{center}
\small
\begin{tabular}{l|l}
\textbf{Layer} & \textbf{Examples} \\ \hline
Application  & \textbf{HTTP} \\
Presentation & \\
Session      & SSL, TLS \\
Transport    & TCP, UDP \\
Network      & IP (naming and routing: DNS, BGP) \\
Data link    & IEEE 802.3 Ethernet (naming and routing: ARP) \\
Physical     & modulation and coding
\end{tabular}
\end{center}
This lecture focuses on the \textbf{top layers}.
\end{defnbox}

\begin{defnbox}[title={HTTP/1.1: HyperText Transfer Protocol}]
The protocol that determines what actions web servers and browsers take in response to commands.
\begin{itemize}
  \item \textbf{Request-response}: (1) the client (browser) sends a \textbf{request} (for an HTML
        file, to update a database, to send a mail, \ldots); (2) the server processes it, performs
        the action, and sends a \textbf{response} (the content to render, or an acknowledgement).
  \item \textbf{Stateless}: each command is executed independently, without any knowledge of the
        previous ones. Neither browser nor server keeps state: \textbf{the state must be sent
        inside the requests and responses}. Very important for security.
\end{itemize}
\end{defnbox}

\begin{defnbox}[title=Cookies]
A \textbf{small piece of data stored by the browser} on the user's device.
\begin{itemize}
  \item \textbf{Main goal}: store state information (shopping cart, login) to create HTTP
        \textbf{sessions} on top of a stateless protocol.
  \item \textbf{Secondary use}: tracking users.
  \item \textbf{Ambient authority in cookies}: if you are logged into \texttt{bank.com}, you have
        cookies for \texttt{bank.com}, and \textbf{every} new HTTP request to \texttt{bank.com}
        automatically includes all of them, so that you stay logged in (no login for every
        request). Whoever makes your browser send a request to \texttt{bank.com} uses your
        session.
\end{itemize}
\end{defnbox}
